\section{Introduction}

Visual anomaly detection is a core component of automated industrial
inspection. A detector must answer two questions: whether a product image is
abnormal, and where the defect is located. Unlike ordinary image
classification, the anomalous class is not available during training and its
appearance can change from one defect to another. The realistic learning
problem is therefore normal-only: model the distribution of defect-free
appearance and flag samples that deviate from it.

Recent methods address this problem in two broad ways. Reconstruction and
synthesis methods learn to reproduce normal appearance or generate pseudo
defects \cite{zavrtanik2021draem,li2021cutpaste}. Patch-distribution methods
estimate compact normal statistics and compare local test features against
them \cite{defard2021padim,roth2022patchcore}. Foundation-model approaches
reuse frozen visual representations learned from much larger datasets. The
last family is especially attractive in a few-shot setting because no
task-specific training is required and the same encoder can be applied to
different products.

AnomalyDINO combines DINOv2 patch tokens with nearest-neighbor retrieval
\cite{damm2025anomalydino,oquab2024dinov2}. Normal reference images populate a
memory bank. Every patch of a test image is assigned the distance to its
closest normal patch, and an image-level score is obtained from the strongest
patch responses. This design naturally provides both a scalar decision and a
spatial anomaly map. The remaining difficulty is that last-layer semantic
features are relatively coarse: they separate obvious semantic anomalies but
can blur thin scratches, small missing parts, and low-contrast contamination.

This project makes three contributions. First, it implements the required
template-difference baseline and exposes how pose, illumination, and
background changes dominate raw pixel differences. Second, it builds a
reproducible AnomalyDINO pipeline on the six categories specified by the course
project, with calibration data excluded from every few-shot reference pool.
Third, it implements a controlled multi-layer feature-fusion variant and
evaluates it with identical data, references, preprocessing, input resolution,
and seeds. We report image AUROC, image AP/F1, patch-level pixel AUROC/AP/F1,
AUPRO, latency, and peak GPU memory.

The full project also studies reliability. Rather than selecting a threshold
from test labels, each category reserves the first 20 normal training images
for calibration and uses their 95th percentile score as the demo OK/NG
threshold. Test labels are touched only after all predictions have been made.
This protocol is stricter than reporting AUROC alone and makes the demo
decision reproducible.
